% Figure from MATH 590 notes, section 18. Compile with the notes preamble.
\begin{tikzpicture}[scale=1.1]
  % open set U
  \draw[black!55, dashed, thick] plot[smooth cycle, tension=0.8] coordinates
    {(-2.3,0.2) (-1.6,1.9) (0.2,2.1) (2.0,1.6) (2.4,-0.3) (1.4,-2.0) (-0.6,-2.1) (-2.1,-1.3)};
  \node[font=\small, black!65] at (2.25,1.75) {$U$};
  % B(x, eps)
  \draw[black!60, dashed] (0,0) circle (1.6);
  \draw[black!60, -stealth] (0,0) -- (-1.13,-1.13);
  \node[font=\scriptsize, black!70] at (-0.72,-0.42) {$\varepsilon$};
  % B(d_i, 1/n)
  \filldraw[figR, fill=figR!15, thick] (0.45,0.25) circle (0.72);
  \draw[figR, -stealth] (0.45,0.25) -- (1.17,0.25);
  \node[font=\scriptsize, figR, above] at (0.87,0.25) {$\tfrac1n$};
  \draw[figR, thin] (0.98,0.74) -- (2.75,1.05);
  \node[font=\scriptsize, figR, right] at (2.75,1.05) {$B(d_i,\tfrac1n)$};
  % countable dense set D (a few of its points)
  \foreach \p in {(-1.9,-0.4),(-1.2,1.3),(-0.3,1.1),(0.9,1.75),(1.9,0.6),(1.3,-1.3),(-0.4,-1.6),(-1.5,-1.4),(0.2,-0.9),(-1.0,0.4),(1.8,-0.5),(-0.9,-0.9)}
    \fill[figB!70] \p circle (0.9pt);
  \fill[figB] (0.45,0.25) circle (1.5pt);
  \node[font=\scriptsize, figB, above left=-1pt] at (0.45,0.25) {$d_i$};
  \fill (0,0) circle (1.5pt);
  \node[font=\scriptsize, below=1pt] at (0,0) {$x$};
  \node[font=\scriptsize, figB] at (-1.35,1.62) {$D$};
\end{tikzpicture}
